% Tabla A1-T7. Generada automáticamente; no editar a mano.
% Requiere \usepackage{booktabs} (y graphicx si la tabla es ancha).
\begin{table}[htbp]
    \centering
    \caption{Análisis I: probabilidad predicha de cada régimen en los percentiles 20, 50 y 80 de la masa estelar y de la metalicidad.}
    \label{tab:A1-T7}
    \small
    \resizebox{\textwidth}{!}{%
    \begin{tabular}{llrrrrrr}
    \toprule
    Variable & Percentil & Valor & Terrestre & Super-Tierra & Tipo Neptuno & Gigante gaseoso & Mayor que Júpiter \\
    \midrule
    $M_\star$ & P20 & 0.777 & 0.122 & 0.300 & 0.463 & 0.114 & 0.002 \\
    $M_\star$ & P50 & 0.945 & 0.101 & 0.279 & 0.455 & 0.158 & 0.007 \\
    $M_\star$ & P80 & 1.110 & 0.081 & 0.251 & 0.432 & 0.210 & 0.025 \\
    {[Fe/H]} & P20 & -0.100 & 0.126 & 0.318 & 0.442 & 0.106 & 0.009 \\
    {[Fe/H]} & P50 & 0.020 & 0.101 & 0.279 & 0.455 & 0.158 & 0.007 \\
    {[Fe/H]} & P80 & 0.144 & 0.077 & 0.235 & 0.452 & 0.230 & 0.005 \\
    \bottomrule
    \end{tabular}
    }
    \par\smallskip
    {\footnotesize\textit{Nota:} Modelo B1; las demás variables se fijan en su mediana.}
\end{table}
